% Part: incompleteness
% Chapter: arithmetization-syntax
% Section: coding-symbols

\documentclass[../../../include/open-logic-section]{subfiles}

\begin{document}

\olfileid{inc}{art}{cod}
\olsection{رمزگذاری نمادها}

زبان پایهٔ~$\Lang L$ از منطق مرتبه‌اول نمادهای زیر را به کار می‌برد:
\[
\lfalse \quad \lnot \quad \lor \quad \land \quad \lif \quad \lforall
\quad \lexists \quad \eq \quad ( \quad ) \quad ,
\]
همراه با مجموعه‌های !!{enumerable} متغیرها و !!{constant}s، و مجموعه‌های
!!{enumerable} !!{function}s و !!{predicate}sیی با آریتی دلخواه. می‌توانیم
به هر یک از این نمادها \emph{رمز}ی چنان نسبت دهیم که هر نماد عددی
یکتا به‌عنوان رمزش بگیرد و هیچ دو نماد متفاوتی عددی یکسان نگیرند.
می‌دانیم این ممکن است زیرا مجموعهٔ همهٔ نمادها !!{enumerable} است و
بنابراین !!a{bijection} میان آن و مجموعهٔ اعداد طبیعی وجود دارد.
اما می‌خواهیم مطمئن شویم نماد، و نیز برخی اطلاعات دربارهٔ آن مانند
آریتیِ یک !!a{function}، را بتوان به‌طور محاسبه‌پذیر از رمزش بازیابی کرد.
البته راه‌های ممکن بسیاری برای این کار هست. اینجا یک راه را می‌آوریم
که از توابع بازگشتی اولیه بهره می‌گیرد. (به یاد آورید
$\tuple{n_0,\dots,n_k}$ عددی است که دنبالهٔ اعداد
$n_0$،~\dots، $n_k$ را رمزگذاری می‌کند.)

\begin{defn}
اگر $s$ نمادی از~$\Lang L$ باشد، \emph{رمز نمادِ}~$\scode s$ را چنین
تعریف کنید:
\begin{enumerate}
\item اگر $s$ در میان نمادهای منطقی باشد، $\scode s$ را جدول زیر
  می‌دهد:
\[
\begin{array}{cccccccccc}
  \lfalse & \lnot & \lor & \land & \lif & \lforall \\
\tuple{0, 0} & \tuple{0, 1} & \tuple{0, 2} & \tuple{0, 3} &
\tuple{0, 4} & \tuple{0, 5} \\
\lexists & \eq & ( & ) & ,\\
\tuple{0, 6} & \tuple{0, 7} &
\tuple{0, 8} & \tuple{0, 9} & \tuple{0, 10}
\end{array}
\]
\item اگر $s$ متغیر $i$اُمِ~$\Obj v_i$ باشد،
$\scode s=\tuple{1,i}$.
\item اگر $s$ !!{constant} $i$اُمِ~$\Obj c_i$ باشد،
  $\scode s=\tuple{2,i}$.
\item اگر $s$ !!{function} $n$موضعیِ $i$اُمِ~$\Obj f_i^n$ باشد،
  $\scode s=\tuple{3,n,i}$.
\item اگر $s$ !!{predicate} $n$موضعیِ $i$اُمِ~$\Obj P_i^n$ باشد،
  $\scode s=\tuple{4,n,i}$.
\end{enumerate}
\end{defn}

\begin{prop}
روابط زیر بازگشتی اولیه‌اند:
\begin{enumerate}
\item $\fn{Fn}(x,n)$ اگر و تنها اگر $x$ برای بعضی~$i$ رمز
$\Obj f_i^n$ باشد؛ یعنی $x$ رمز یک !!{function} $n$موضعی باشد.
\item $\fn{Pred}(x,n)$ اگر و تنها اگر $x$ برای بعضی~$i$ رمز
$\Obj P_i^n$ باشد، یا $x$ رمز~$\eq$ و~$n=2$ باشد؛ یعنی $x$ رمز یک
!!{predicate} $n$موضعی باشد.
\end{enumerate}
\end{prop}

\begin{defn}
اگر $s_0$،~\dots، $s_{n-1}$ دنباله‌ای از نمادها باشد، \emph{عدد
گودلِ} آن برابر است با
$\tuple{\scode{s_0},\dots,\scode{s_{n-1}}}$.
\end{defn}

\begin{explain}
توجه کنید \emph{رمزها} و \emph{اعداد گودل} چیزهایی متفاوت‌اند. برای
نمونه، متغیر~$\Obj v_5$ رمز
$\scode{\Obj v_5}=\tuple{1,5}=2^2\cdot3^6$ را دارد. اما متغیر
$\Obj v_5$ وقتی به‌عنوان ترم در نظر گرفته شود دنباله‌ای از نمادها
(به طول~$1$) نیز هست. \emph{عدد گودلِ}~$\Gn{\Obj v_5}$ برای
\emph{ترمِ}~$\Obj v_5$ برابر است با
$\tuple{\scode{\Obj v_5}}=2^{\scode{\Obj v_5}+1}
=2^{2^2\cdot3^6+1}$.
\end{explain}

\begin{ex}
به یاد آورید اگر $k_0$،~\dots، $k_{n-1}$ دنباله‌ای از اعداد باشد،
رمز دنبالهٔ~$\tuple{k_0,\dots,k_{n-1}}$ در رمزگذاریِ توان‌های اعداد
اول برابر است با
\[
2^{k_0+1}\cdot3^{k_1+1}\cdot \dots \cdot p_{n-1}^{k_{n-1}+1},
\]
که در آن $p_i$ عدد اول $i$اُم است (با آغاز از~$p_0=2$). پس برای
نمونه، فرمول~$\eq[\Obj v_0][\Obj 0]$ یا، به‌طور صریح‌تر،
${\eq}(\Obj v_0,\Obj c_0)$ عدد گودل زیر را دارد:
\[
\tuple{\scode{\eq},\scode{(},\scode{\Obj v_0},\scode{,}, \scode{\Obj
    c_0},\scode{)}}.
\]
اینجا $\scode{\eq}$ برابر
$\tuple{0,7}=2^{0+1}\cdot3^{7+1}$، $\scode{\Obj v_0}$ برابر
$\tuple{1,0}=2^{1+1}\cdot3^{0+1}$ و به همین ترتیب است. پس
$\Gn{=(\Obj v_0,\Obj c_0)}$ برابر است با
\begin{multline*}
2^{\scode{=} + 1}\cdot 3^{\scode{(}+1}\cdot 5^{\scode{\Obj v_0}+1}
\cdot 7^{\scode{,} + 1} \cdot 11^{\scode{\Obj c_0}+1} \cdot
13^{\scode{)}+1} = \\
2^{2^1\cdot 3^8 + 1}\cdot 3^{2^1\cdot 3^9+1}\cdot 5^{2^2\cdot 3^1+1}
\cdot 7^{2^1\cdot 3^{11} + 1} \cdot 11^{2^3\cdot3^1+1} \cdot
13^{2^1\cdot3^{10}+1} = \\
2^{13\,123}\cdot 3^{39\,367}\cdot 5^{13}\cdot 7^{354\,295}\cdot11^{25}\cdot13^{118\,099}.
\end{multline*}
\end{ex}

\end{document}
